% =============================================================================
%  CONCLUSION GÉNÉRALE ET PERSPECTIVES
% =============================================================================
\cleardoublepage
\entete{Conclusion générale}
\chapter*{CONCLUSION GÉNÉRALE}
\addcontentsline{toc}{chapter}{CONCLUSION GÉNÉRALE}

Au terme de ce travail, mené du 20 juin au 20 août 2026 au sein du \emph{Think
Tank} Fou'ndjem de la Fondation Jean Félicien Gacha, le point de départ mérite
d'être rappelé : la découverte de la Fondation supposait de s'y rendre, ce qui
écartait les publics éloignés, la diaspora et les personnes à mobilité réduite,
et là où une présence numérique existait, elle se réduisait à une vitrine figée.
La question posée était donc celle d'une plateforme unifiée capable d'offrir une
visite en trois dimensions, de retrouver les œuvres apparentées dans d'autres
institutions et d'assurer une assistance fiable fondée sur les seules données de
la Fondation.

Pour y répondre, nous avons conduit une recherche interactive de niveau
intégrateur, selon une approche mixte associant l'analyse documentaire,
l'observation directe, six entretiens semi-directifs et une enquête auprès de
quarante-deux visiteurs. Le travail a abouti à une plateforme opérationnelle,
MUSÉA, articulant vingt fonctionnalités autour de quatre ensembles : un
progiciel de gestion couplé à un site public éditable ; une visite immersive à
360~degrés complétée par la 3D et la réalité augmentée à taille réelle ; un
assistant conversationnel ancré sur les documents validés ; et un dispositif de
rapprochement avec les collections ouvertes du monde. Les quatre hypothèses
spécifiques ont été vérifiées et l'hypothèse générale corroborée.

Au-delà de l'outil livré, trois enseignements constituent la contribution propre
de ce travail.

\begin{enumerate}
  \item \textbf{Le décalage de vocabulaire constitue une seconde dispersion des
        œuvres.} Une recherche en français dans les catalogues internationaux
        ramène du bruit là où la même recherche, reformulée dans leur
        vocabulaire de catalogage, identifie des correspondances précises : le
        meilleur score passe de 57 à 85 sur 100. À la dispersion physique
        s'ajoute une dispersion documentaire, réversible celle-là.

  \item \textbf{L'isolement des données ne peut pas être une propriété du
        code.} Les deux fuites détectées étaient d'origine applicative : seule
        une garantie portée par la base de données résiste à l'erreur de
        programmation.

  \item \textbf{La qualité d'une visite restituée repose sur quatre
        continuités}, spatiale, d'échelle, documentaire et narrative, dont la
        rupture d'une seule fait échouer l'ensemble.
\end{enumerate}

Sur le plan personnel, ce stage aura permis de conduire un projet complet, du
besoin exprimé jusqu'à la mise en production. Mené auprès d'une institution
unique, ce travail ne prétend pas à la généralisation : sa transposition
demandera vérification, objet des perspectives qui suivent.

% =============================================================================
\cleardoublepage
\entete{Perspectives}
\chapter*{PERSPECTIVES}
\addcontentsline{toc}{chapter}{PERSPECTIVES}

Le dispositif construit ouvre un ensemble de prolongements, que nous présentons
regroupés en quatre axes, du plus immédiat au plus prospectif. Certains d'entre
eux n'ont pu être menés à terme durant le stage en raison de contraintes de
temps, de moyens ou d'accès ; nous le signalons explicitement le cas échéant.

\section*{Axe 1 : Consolider ce qui existe}
\addcontentsline{toc}{section}{Axe 1 : Consolider ce qui existe}

\begin{enumerate}
  \item \textbf{Comparer les œuvres sur l'image et non sur le texte.} C'est la
        perspective la plus importante. Les rapprochements documentaires
        reposent actuellement sur la description écrite des œuvres, faute de
        processeur graphique et de possibilité d'installer un modèle de vision
        sur le poste de travail. Le recours à un modèle de plongement visuel
        permettrait de rapprocher deux objets visuellement proches mais décrits
        selon des conventions différentes, situation fréquente entre un
        catalogue africain et un catalogue européen. Cette évolution supprimerait
        la principale limite de la mémoire réunifiée.

  \item \textbf{Produire les cartes de profondeur en production.} Le moteur de
        relief a été vérifié sur une carte fabriquée pour l'essai, mais aucune
        carte n'a encore été produite sur des œuvres réelles : le traitement
        exige des bibliothèques qui n'ont pu être installées. Une exécution dans
        un conteneur, telle que prévue par l'intervention proposée, lèverait cet
        obstacle.

  \item \textbf{Élargir l'angle de parallaxe du relief.} L'illusion est
        actuellement bornée à vingt degrés, au-delà desquels les zones absentes
        de la photographie s'étirent. Une passe de reconstitution des zones
        masquées permettrait d'élargir cet angle.

  \item \textbf{Obtenir les clés d'accès aux quatre collections écartées} : Europeana, le Rijksmuseum, la Smithsonian Institution et les Harvard Art
        Museums. Leur intégration est déjà prévue techniquement ; elle ne dépend
        que d'une démarche administrative.

  \item \textbf{Éprouver la chaîne de rapprochement de bout en bout}, chaque
        maillon n'ayant été vérifié qu'isolément, et \textbf{étendre les essais
        de réalité augmentée} à un parc d'appareils représentatif.
\end{enumerate}

\section*{Axe 2 : Enrichir l'expérience du visiteur}
\addcontentsline{toc}{section}{Axe 2 : Enrichir l'expérience du visiteur}

\begin{enumerate}
  \item \textbf{Le mur de la dispersion.} L'écran de fin de parcours existe ;
        il reste à y animer le globe pour annoncer au visiteur : « les œuvres
        apparentées à celles que vous venez de voir se trouvent dans tel nombre
        de pays ». Le geste narratif serait fort : la visite s'achèverait sur la
        mesure de ce dont l'institution a été séparée.

  \item \textbf{Un guide vocal contextuel.} L'assistant sait déjà où se trouve
        le visiteur. Il pourrait commenter spontanément la salle où celui-ci
        vient d'entrer, à la manière d'un guide qui prend la parole sans qu'on
        l'interroge.

  \item \textbf{Le mode hors connexion.} L'application est déjà installable ; il
        resterait à mettre en cache les panoramas et les modèles des salles
        visitées, afin qu'une visite entamée puisse se poursuivre sur une
        connexion défaillante, situation courante pour une partie du public
        visé.

  \item \textbf{La reconnaissance d'une œuvre par photographie}, pour qu'un
        visiteur présent sur place obtienne immédiatement la notice d'une pièce
        sans la chercher dans un catalogue.

  \item \textbf{La médiation ludique à destination des scolaires}, que les
        mécanismes de parcours à étapes déjà présents dans la base
        permettraient de soutenir, en présence comme à distance.

  \item \textbf{Les casques de réalité virtuelle}, écartés à dessein pour des
        raisons d'accessibilité, mais envisageables en complément pour un usage
        sur place ou en médiathèque.
\end{enumerate}

\section*{Axe 3 : Ancrer le dispositif dans la culture locale}
\addcontentsline{toc}{section}{Axe 3 : Ancrer le dispositif dans la culture locale}

\begin{enumerate}
  \item \textbf{Le multilinguisme étendu aux langues du territoire.} L'ajout des
        langues de l'Ouest camerounais (ghomálá', fe'efe'e, medumba) à
        l'application, aujourd'hui bilingue, transformerait la portée du projet :
        il ne s'agirait plus seulement de faire connaître un patrimoine à
        l'extérieur, mais de le restituer aux communautés qui en sont issues,
        dans leur langue. La structure du système le permet sans modification.

  \item \textbf{La collecte des récits oraux.} Le patrimoine immatériel reste
        peu documenté. Une campagne d'enregistrement auprès des détenteurs de la
        mémoire, rattachée aux fiches d'œuvres et aux personnages de la
        généalogie, enrichirait la base documentaire de l'assistant.

  \item \textbf{L'approfondissement de la généalogie}, qui gagnerait à intégrer
        les migrations historiques, les alliances entre chefferies et les corpus
        d'objets rattachés à chaque règne.
\end{enumerate}

\section*{Axe 4 : Industrialiser, pérenniser et essaimer}
\addcontentsline{toc}{section}{Axe 4 : Industrialiser, pérenniser et essaimer}

\begin{enumerate}
  \item \textbf{Le passage aux conteneurs orchestrés}, décrit au chapitre~4,
        préalable à l'automatisation de l'ingestion documentaire ; puis, au-delà
        d'une dizaine d'institutions, la migration vers Kubernetes, que la
        conteneurisation rend possible sans réécriture.

  \item \textbf{Le suivi de la dépense et la sauvegarde.} L'étiquetage des
        ressources par institution fonderait une tarification équitable, et une
        sauvegarde externalisée devrait être formalisée, les données
        patrimoniales étant irremplaçables.

  \item \textbf{L'évaluation d'impact}, origine des visiteurs, durée des
        parcours, taux d'adhésion, retombées sur la fréquentation physique.

  \item \textbf{L'accueil d'institutions partenaires} et \textbf{l'ouverture des
        données de la Fondation} selon les mêmes standards que les catalogues
        qu'elle interroge, afin que le mouvement devienne réciproque et que les
        rapprochements validés alimentent le dialogue avec les institutions
        détentrices.

  \item \textbf{Le modèle économique de la plateforme.} La mutualisation du
        catalogue mondial fait décroître le coût par institution à mesure que le
        réseau s'étoffe : de petites structures accèdent ainsi à un outil
        qu'elles ne pourraient financer isolément.
\end{enumerate}
